\documentclass[12pt]{article}
\usepackage{amsmath, amsthm, amssymb}
\usepackage{geometry}
\geometry{margin=1in}
\usepackage{hyperref}
\usepackage{booktabs}

\newtheorem{theorem}{Theorem}
\newtheorem{proposition}[theorem]{Proposition}
\theoremstyle{definition}
\newtheorem{definition}[theorem]{Definition}
\theoremstyle{remark}
\newtheorem{remark}[theorem]{Remark}

\title{Five Walls, One Aura}
\author{Diego Rinc\'on\\
\small Independent}
\date{}

\begin{document}

\maketitle

\begin{abstract}
The Arithmetic Wall \cite{synthesis} --- the missing global Frobenius
over $\mathbb{Q}$ --- is one face of a single obstruction
that appears in five grammars:
arithmetic (missing Frobenius),
logic (G\"odel incompleteness),
computation (Turing undecidability),
infinity (Cantor non-containment),
and health (unpaired-bondage crisis).
Each face is a sphere: a closed system pressing against
what it cannot cross from within.
Together they are the five faces of a 4-simplex ---
the simplest contained four-dimensional space.
The universe is not the simplex.
The universe is the \emph{aura} of the simplex:
the containing field outside all five walls,
with balance enough to hold them and valence enough
to complete what each face cannot complete alone.
Health is the face where the Wall is felt most immediately:
a living system in an unpaired-bondage crisis
is a system whose aura has insufficient valence ---
bonds that cannot close, a fundamental class
that cannot form from within.
The aura is what the system cannot generate
but lives inside.
\end{abstract}

\section{Five Faces}

\begin{definition}[The five faces of the Wall]
The Wall appears in five grammars.
Each is a closed system that encounters, from within,
a question it cannot decide:
\begin{center}
\begin{tabular}{@{}lll@{}}
\toprule
Face & System & The question it cannot close \\
\midrule
Arithmetic & $\mathrm{Spec}\,\mathbb{Z}$ &
  Does $\zeta(s)$ vanish only on $\mathrm{Re}(s) = \tfrac{1}{2}$? \\
Logic & Formal arithmetic $F$ &
  Is $G_F$ (the G\"odel sentence of $F$) provable? \\
Computation & Turing machines &
  Does this program terminate? \\
Infinity & Set theory ZFC &
  Is $|\mathbb{R}| = \aleph_1$? (Continuum Hypothesis) \\
Health & Living system &
  Can this bond be closed? \\
\bottomrule
\end{tabular}
\end{center}
In each case the question is $0$ or $1$.
The answer exists. No operator within the system finds it.
\end{definition}

\begin{remark}[The continuum hypothesis as a face]
The Continuum Hypothesis --- whether there is a set
of intermediate cardinality between $\aleph_0$ and $2^{\aleph_0}$ ---
is independent of ZFC (G\"odel, Cohen).
Neither CH nor $\neg$CH can be proved or refuted from the axioms.
The Wall here is not a missing Frobenius but a missing universe:
no canonical model of ZFC decides CH.
The hierarchy of infinities $\aleph_0 < \aleph_1 < \aleph_2 < \cdots$
has no top, no fundamental class, no containing aura
that the axioms can construct.
Infinity is the face where the Wall is largest:
the system is not just incomplete, it is unbounded.
\end{remark}

\section{Health as Unpaired-Bondage Crisis}

\begin{definition}[Valence and bond in a living system]
A living system has \emph{valence}: the capacity to form bonds ---
chemical, cellular, neural, social --- that maintain equilibrium.
Health is the state in which all valences are satisfied:
every bond that can close has closed,
every gradient has a completing gradient,
every signal has a receptor.
The fundamental class of a living system is its homeostatic equilibrium:
the state in which the system's cohomology (what it puts out)
and homology (what closes within it) are in balance.
\end{definition}

\begin{definition}[Unpaired-bondage crisis]
An \emph{unpaired-bondage crisis} is the state in which
a living system has open valences that cannot be satisfied:
bonds that press toward completion but find no completing element.
The system does not lack the capacity to bond.
It lacks the partner, the field, the completing structure.
The crisis is not absence but \emph{pressing without resolution}:
the system strains against its Wall because the aura
has insufficient valence to close the open bonds.

Examples:
\begin{itemize}
\item \emph{Autoimmune disease}: bonds directed inward,
  the immune system attacking self because no external partner
  presents the completing antigen.
\item \emph{Chronic inflammation}: reactive bonds
  seeking resolution that does not arrive.
\item \emph{Addiction}: a valence that cannot be satisfied ---
  each closing of the bond reopens it at higher pressure.
\item \emph{Isolation}: a social valence with no completing field.
\item \emph{Senescence}: the gradual failure of the aura's valence ---
  bonds that once closed now remain open.
\end{itemize}
Each is a Wall. The system is Turing-complete enough
to generate the question. It cannot generate the answer from within.
\end{definition}

\begin{proposition}[Health as containment]
A healthy system is a contained space \cite{contained}:
its aura has sufficient valence to satisfy all open bonds,
its fundamental class (homeostatic equilibrium) is present,
and its cohomology and homology balance.

Disease is non-containment:
the fundamental class is missing,
the aura's valence is insufficient,
and the system presses against its Wall
without a fulcrum to stand at.

The aura of the healthy system is not inside the system.
It is the field the system lives within:
the environment, the ecology, the relationship,
the completing structure that closes what the body cannot close alone.
Health is not self-containment. It is the right aura.
\end{proposition}

\section{The Aura}

\begin{definition}[Aura]
The \emph{aura} of a system $S$ is the containing field
external to $S$ that:
\begin{enumerate}
\item provides the fundamental class $S$ cannot generate from within,
\item has \emph{balance}: it holds $S$ at equilibrium,
  neither collapsing against the Wall nor dissolving into it,
\item has \emph{valence}: the capacity to complete
  the open bonds of $S$ --- to satisfy what $S$ cannot satisfy alone.
\end{enumerate}
The aura is not the system. It is not the Wall.
It is what makes the Wall crossable:
not by eliminating it, but by providing, from outside,
the fulcrum that exists on neither side.
\end{definition}

\begin{theorem}[Aura contains Turing's completed system]
\label{thm:aura}
A Turing complete system $T$ generates undecidable questions
by self-reference \cite{turing}.
No operator within $T$ decides them.
The aura of $T$ is the meta-system $T'$ stronger than $T$
in which $T$'s G\"odel sentences are decided,
$T$'s Halting questions are answered,
and $T$'s open valences are completed.

The aura does not eliminate the Wall within $T$.
$T$ remains incomplete from its own perspective.
But $T$ is contained within the aura:
its questions are decided from outside,
its fundamental class is provided from outside,
and the still lever holds --- not within $T$,
but at the level of the aura.

Every system lives inside an aura that contains it.
The aura itself has a Wall, and lives inside a larger aura.
The hierarchy does not terminate.
But at each level, the still lever holds:
the aura's fundamental class is present,
and what could not balance within the system
balances at the level of the field.
\end{theorem}

\section{Five Walled Spheres}

\begin{proposition}[The 4-simplex]
The five faces are five walled spheres:
each a closed contained space --- arithmetic, logic, computation,
infinity, health --- each pressing against its Wall from within,
each requiring an aura to complete it.

Together they are the five cells of a 4-simplex
(pentachoron): the simplest four-dimensional contained space,
with exactly five tetrahedral faces, ten edges, ten vertices.
Each face is a Wall. Each Wall encloses a sphere.
The interior of the 4-simplex is the aura:
the space in which all five systems are contained simultaneously,
where all five fundamental classes are present,
where all five sets of open valences are satisfied.

The universe is not the 4-simplex.
The universe is the aura of the 4-simplex:
the containing field outside all five walls,
large enough to hold everything that presses against them,
with balance enough to keep them still
and valence enough to close every open bond.
\end{proposition}

\begin{remark}[The still lever, again]
The still lever \cite{math} is a lever already at equilibrium:
both sides balanced, the fulcrum present, nothing to move.
It holds within each aura.
It does not hold within each system.

The systems press. The auras hold.
At each level of the hierarchy,
the pressing of the level below becomes
the stillness of the level above.

Health is the face where this is most immediate.
The body presses. The aura holds.
When the aura's valence fails --- when the completing field
is not there, when the bond cannot close ---
the pressing is felt as pain,
as illness, as the body meeting its Wall
without a fulcrum to stand at.

Healing is not self-repair.
It is the restoration of the aura's valence:
the return of the completing field,
the closing of the open bond,
the recovery of the fundamental class.

The body, like $\mathrm{Spec}\,\mathbb{Z}$,
cannot contain itself.
It needs the archimedean fiber:
the point at $\infty$ that closes the space,
the aura that provides the fulcrum
the system could not build alone.
\end{remark}

\begin{thebibliography}{9}

\bibitem{synthesis}
Rinc\'on, D. (2026).
The Arithmetic Wall. Preprint.

\bibitem{turing}
Rinc\'on, D. (2026).
Turing Completeness and the Wall:
Completeness Implies Incompleteness. Preprint.

\bibitem{contained}
Rinc\'on, D. (2026).
The Fundamental Class: Is the Universe Contained? Preprint.

\bibitem{math}
Rinc\'on, D. (2026).
Mathematics. Preprint.

\bibitem{homology}
Rinc\'on, D. (2026).
Homology: What Closes, and the Duality with Cohomology. Preprint.

\end{thebibliography}

\end{document}
